% Inactive proof fragment; uses the existing book macros and dependencies.
\subsection{Left limits of greatest exponential solutions}

Keep the fixed compact connected Hermitian manifold $(X,h)$, nonzero nonnegative density $f\in L^1(\mathrm dV)$ and greatest solutions $U_\lambda$ for $\lambda\in\Lambda_h(f)$. The corrected order in \cref{lem:positive-exponent-uniform-upper-barrier} also controls approach from the left. Its limit solves the limiting equation, although greatest maximality alone does not identify it with the greatest solution there.

\begin{proposition}\label{prop:greatest-exponential-left-limit}
Suppose $\Lambda_h(f)\ne\varnothing$, and put $a_0=\inf\Lambda_h(f)$. For every $a>a_0$, there is a genuine $L_a\in\PSH(X,h)$ such that
\[
 U_\lambda\longrightarrow L_a
       \quad\text{in }L^1\text{ and in capacity}\quad(\lambda\uparrow a).
\]
It satisfies
\[
 L_a\leq U_a,\qquad
 \mathrm{MA}^{\mathrm{ext}}_h(L_a)=\mathrm e^{aL_a}f\,\mathrm dV.
\]
The equation measures converge in total variation to this finite measure. The left limit is independent of the auxiliary exponent and of the cofinal sequence used to construct it.
\end{proposition}
\begin{proof}
Choose $\tau\in\Lambda_h(f)$ with $0<\tau<a$; the upper-interval property gives such a choice. Put $K=K_\tau$ from \cref{lem:positive-exponent-uniform-upper-barrier} and
\[
 V_\lambda=U_\lambda+K\log(\lambda/\tau)
       \qquad(\lambda\in\Lambda_h(f),\ \lambda\geq\tau).
\]
This family is increasing. For a sequence $\lambda_j\uparrow a$ with $\tau\leq\lambda_j<a$, its members lie between the genuine integrable potentials $V_\tau$ and $V_a$. The upper-envelope construction in \cref{sec:obstacle-envelopes} gives
\[
 V^-_a=\left(\sup_{\tau\leq\lambda<a}V_\lambda\right)^*
        =\left(\lim_jV_{\lambda_j}\right)^*\in\PSH(X,h).
\]
Indeed cofinality and the order give the same pointwise supremum for every such sequence. Regularization changes it only on a pluripolar set. Dominated convergence between the integrable bounds gives $L^1$ convergence to $V^-_a$, and \cref{lem:monotone-psh-capacity} gives convergence in capacity. Since the added constants converge, set
\[
 L_a=V^-_a-K\log(a/\tau).
\]
Then $U_{\lambda_j}\to L_a$ in both topologies and $L_a\leq U_a$; the latter inequality survives regularization because $V_a$ is upper semicontinuous.

The potentials $U_{\lambda_j}-K$ and their limit $L_a-K$ are nonpositive. Apply \cref{thm:exponential-closedness} with exponents tending to $a>0$ and densities $\mathrm e^{\lambda_jK}f\to\mathrm e^{aK}f$ in $L^1$. It proves the exact equation and total-variation convergence.

Every sequence approaching $a$ from the left has a cofinal increasing subsequence. If potential convergence in either topology, or total-variation convergence of the equation measures, failed, a sequence staying outside a fixed neighbourhood of the proposed limit would therefore have an increasing subsequence, contrary to the preceding argument. This proves the full left limit and measure convergence. Repeating the construction with another $\tau$ gives the same $L^1$ limit; canonical quasi-plurisubharmonic representatives that agree almost everywhere agree everywhere. Thus the auxiliary choice does not change $L_a$.
\end{proof}

\begin{proposition}\label{prop:greatest-exponential-countable-jumps}
On $(a_0,+\infty)$, the set
\[
 D=\{a>a_0:L_a\ne U_a\}
\]
is at most countable. At every $a\notin D$, the greatest solutions are continuous in $L^1$ and in capacity as $\lambda\to a$, and their equation measures converge in total variation to $\mathrm e^{aU_a}f\,\mathrm dV$.
\end{proposition}
\begin{proof}
Fix $a_0<\alpha<\beta$, choose an admitted $0<\tau<\alpha$, and use the corrected family $V_\lambda$ from the preceding proof on $[\alpha,\beta]$. Its integrals
\[
 I(\lambda)=\int_XV_\lambda\,\mathrm dV
\]
form a finite bounded increasing real function. For $a\in(\alpha,\beta]$, the left-limit proposition gives
\[
 I(a)-I(a-)=\int_X(U_a-L_a)\,\mathrm dV.
\]
If $L_a\ne U_a$, this quantity is positive: the integrable nonnegative difference cannot vanish almost everywhere without equality of the canonical potentials. For any finite ordered collection of such points, the sum of their jumps is at most $I(\beta)-I(\alpha)$. Consequently there are only finitely many jumps at least $1/m$ for each positive integer $m$, and at most countably many positive jumps. A countable cover by intervals with rational endpoints inside $(a_0,+\infty)$ proves the assertion about $D$.

If $a\notin D$, the left limit is $U_a$, and \cref{prop:greatest-exponential-right-continuity} supplies the same right limit in both topologies and for the equation measures. Splitting an arbitrary approaching sequence into its two sides proves the stated two-sided convergence.
\end{proof}

Whether an actual jump occurs, or all left limits equal the greatest solution, remains undecided. The statements concern positive exponents in the interior of the solvable interval; they do not supply a zero-exponent limit or an intrinsic criterion for the interval's lower endpoint.
